% FPU sequencer and the four floating-point datapaths.
% Requires tikz_styles.tex.
\begin{figure}[p]
\centering
\begin{tikzpicture}[x=1cm,y=1cm]
  \node[rf, minimum width=26mm] (ops) at (0,1.3) {binary64\\or complex};
  \node[ctl, minimum width=24mm] (seq) at (3.2,1.3) {sequencer};
  \node[ctl, minimum width=28mm] (kind) at (3.2,2.7) {comb / direct / sequence};
  \node[acc, minimum width=20mm] (add) at (0,0) {fp\_add\\4 cyc};
  \node[acc, minimum width=20mm] (mul) at (2.6,0) {fp\_mul\\6 cyc};
  \node[acc, minimum width=28mm] (div) at (5.6,0) {divrem\\div, fmod, sqrt};
  \node[acc, minimum width=24mm] (pow) at (2.6,-1.5) {fp\_pow};
  \node[acc, minimum width=24mm] (um) at (5.6,-1.5) {75$\times$53 mul\\shared};
  \draw[arr] (ops) -- (seq);
  \draw[arr] (kind) -- (seq);
  \draw[arr] (seq) -- (add);
  \draw[arr] (seq) -- (mul);
  \draw[arr] (seq) -- (div);
  \draw[arr] (mul) -- (pow);
  \draw[arr] (mul.east) -- ++(0.25,0) |- (um.north);
  \draw[arr] (pow) -- (um);
\end{tikzpicture}
\caption{FPU (\texttt{pycore\_fpu}).
One adder, one multiplier, one divider, and one power unit.
The divider is also the square root.
The significand multiplier is a single \texttt{pycore\_umul\_seq} of 75$\times$53, step 14, muxed between \texttt{fp\_mul} (53$\times$53, zero-extended) and \texttt{fp\_pow}.
The sequencer never runs those two at once.
Scalar add, subtract, multiply, and divide pass through.
Remainder, floor division, general powers, and complex arithmetic are sequences of those datapaths in CPython's order.
Compares, \texttt{not}, and unary plus and minus are combinational.}
\label{fig:fpu}
\end{figure}

\begin{figure}[htbp]
\centering
\begin{tikzpicture}[node distance=2.4mm]
  \node[ctl, minimum width=62mm] (c0) {C0 \quad unpack, order by magnitude};
  \node[ctl, minimum width=62mm, below=of c0] (c1) {C1 \quad align, sticky, add or subtract};
  \node[ctl, minimum width=62mm, below=of c1] (c2) {C2 \quad leading-zero normalize};
  \node[acc, minimum width=62mm, below=of c2] (c3) {C3 \quad round to nearest even, pack};
  \foreach \p/\q in {c0/c1,c1/c2,c2/c3} \draw[arr] (\p) -- (\q);
\end{tikzpicture}
\caption{Binary64 add and subtract (\texttt{pycore\_fp\_add}).
Three registered stages; the round and pack of the last stage are combinational, so \texttt{done} is four cycles after accept.
Specials (NaN, infinity) can be decided in C0 and carried forward.
Subnormals, signed zeros, and the guard/round/sticky scheme are in the shared packer \texttt{pycore\_f64\_pack\_round}.
Each stage already has its own valid bit so the unit can later accept a new operation while an old one drains.}
\label{fig:fpadd}
\end{figure}

\begin{figure}[htbp]
\centering
\begin{tikzpicture}[node distance=2.4mm]
  \node[ctl, minimum width=64mm] (c0) {C0 \quad unpack, normalize subnormals};
  \node[acc, minimum width=64mm, below=of c0] (st) {C1--C4 \quad 53$\times$53 on the shared multiplier};
  \node[acc, minimum width=64mm, below=of st] (pk) {C5 \quad normalize one bit, round, pack};
  \foreach \p/\q in {c0/st,st/pk} \draw[arr] (\p) -- (\q);
  \node[anno, anchor=west] at ($(st.east)+(0.25,0)$) {STEP $= 14$\\$\lceil 53/14\rceil = 4$};
\end{tikzpicture}
\caption{Binary64 multiply (\texttt{pycore\_fp\_mul}).
The 53$\times$53 significand product is iterative.
At the default step, \texttt{done} is the sixth cycle counting accept as the first.
The parent owns the multiplier; this unit only presents operands and waits for the product.}
\label{fig:fpmul}
\end{figure}

\begin{figure}[htbp]
\centering
\begin{tikzpicture}[node distance=3mm and 6mm]
  \node[acc, minimum width=34mm] (divu) {restoring core\\radix 4};
  \node[ctl, above=of divu] (div) {divide\\56 quotient bits\\32 cycles};
  \node[ctl, left=of divu] (fmod) {fmod\\exact\\by exponent gap};
  \node[ctl, right=of divu] (sqrt) {sqrt\\56 root bits\\32 cycles};
  \draw[arr] (div) -- (divu);
  \draw[arr] (fmod) -- (divu);
  \draw[arr] (sqrt) -- (divu);
  \node[anno, align=left, anchor=north] at ($(divu.south)+(0,-0.2)$)
    {Specials finish in 3 cycles. $|a|<|b|$ fmod finishes in 4.\\
     Otherwise fmod takes $5+\lceil(d+1)/2\rceil$ cycles, $d$ the exponent difference.};
\end{tikzpicture}
\caption{Divider, remainder, and square root (\texttt{pycore\_fp\_divrem} on \texttt{pycore\_udiv\_seq}).
Division iterates 56 quotient bits, two per cycle.
Square root uses the same remainder register and the same three subtractors, with the divisor multiples replaced by the root-digit candidates $q(8Q+q)$ for $q \in \{1,2,3\}$.
The final remainder is the sticky bit, so the root is correctly rounded.
The sequencer sends $x^{0.5}$ here (35 cycles including the power special-case setup) instead of through the general power unit.}
\label{fig:fpdiv}
\end{figure}

\begin{figure}[p]
\centering
\begin{tikzpicture}[node distance=2.6mm]
  \node[rf, minimum width=70mm] (in) {$x = m'\,2^{E}$, \quad $m'$ in $[2/3,\,4/3)$};
  \node[acc, minimum width=70mm, below=of in] (log)
    {log2 \quad signed-digit normalization \quad ROM of $\log_2(1+d_j 2^{-j})$};
  \node[acc, minimum width=70mm, below=of log] (mul)
    {multiply \quad $L = E + \log_2 m'$ by the 53-bit significand of $y$};
  \node[trap, minimum width=70mm, below=of mul] (lim)
    {$|P| \ge 2^{11}$ \quad overflow to $+\infty$ or underflow to $+0$};
  \node[acc, minimum width=70mm, below=of lim] (exp)
    {exp2 \quad restoring recurrence on the fraction of $P$};
  \node[ctl, minimum width=70mm, below=of exp] (pk)
    {pack \quad $Z \cdot 2^{\lfloor P \rfloor}$, round to nearest even};
  \foreach \p/\q in {in/log,log/mul,mul/lim,lim/exp,exp/pk} \draw[arr] (\p) -- (\q);
\end{tikzpicture}
\caption{Binary64 power (\texttt{pycore\_fp\_pow}): $x^{y} = 2^{y \log_2 x}$ for positive finite $x$ and finite $y$.
Both transcendental steps are digit recurrences on one fixed-point datapath (two barrel shifters and two adders per cycle).
The only multiplier is the shared 75$\times$53 unit.
Log2 is about 80 cycles, the multiply 5, and exp2 about 72, plus a few cycles of setup.
A power-of-two base skips the log stage.
The sequencer adds the CPython special cases around this unit, so a general \texttt{float ** float} is 171 cycles, 91 for a power-of-two base, and 19 when overflow or underflow is decided at the multiply.
$x^{\pm 1}$, $x^{2}$, and $x^{-2}$ stay on multiply and divide instead.}
\label{fig:fppow}
\end{figure}

\begin{figure}[htbp]
\centering
\begin{tikzpicture}[node distance=2.2mm and 3mm]
  \node[acc, minimum width=28mm] (add) {$+$ / $-$\\two adds\\10 cyc};
  \node[acc, right=of add, minimum width=32mm] (mul) {$\times$\\4 muls, 2 adds\\34 cyc};
  \node[acc, right=of mul, minimum width=36mm] (div) {$/$ Smith scaling\\3 div, 3 mul, 3 add\\128 cyc};
  \node[ctl, below=of div] (rec) {inf / zero recovery\\two more adds};
  \draw[arr] (div) -- (rec);
\end{tikzpicture}
\caption{Complex sequences, all on the scalar datapaths.
A complex value is two binary64 lanes.
Addition and subtraction are two adds.
Multiplication is CPython's four-multiply product.
Division is Smith's algorithm.
When that algorithm would return \texttt{nan+nanj} and exactly one operand holds an infinity, two further adds rebuild the result from signs.
Finite operands skip that recovery and keep the 128-cycle path.
A negative base raised to a fractional real power is not a complex result here; it traps \texttt{FPU\_EXCEPTION}.}
\label{fig:cplx}
\end{figure}
